\chapter{Literature Review}\label{ch:litreview}

\begin{draftnote}
MODULE 3 --- first draft, written 2026-08-02 (Day 25, evening).

\textbf{The old ``BLOCKED: no reference PDFs are in this repo'' note is superseded and has been
deleted.} The chapter is written from the 21 verified entries in \texttt{references.bib} against
the spine agreed in \texttt{logbook/10\_references.md}. Every citation here appears in that claim
map; if a new claim is added, add its row to the map first.

\textbf{Claim discipline used.} Load-bearing technical claims are drawn from the constrained-RL
foundations and from the three papers held in full (\cite{khan2026rl_precision_grasping},
\cite{shahid2022continuous_grasping}, \cite{ferreira2025grasping}). The surveys and application
papers are cited for framing, taxonomy and positioning only, and no specific numerical result is
attributed to a paper not held. Section~\ref{sec:b-gap} is load-bearing --- it is where the gap
statement is made, and Chapter~\ref{ch:intro} should restate it.

\textbf{Open for revision:} target is 8--10 pages, check against the built PDF. Consider whether
Section~\ref{sec:b-repro} belongs here or in Chapter~\ref{ch:methods} --- it is placed here
because it motivates the ten-seed protocol, but it is arguably method rather than background.
\end{draftnote}

\section{The safety requirement in learned manipulation}\label{sec:b-motivation}

A robot manipulator that learns its own control policy poses a problem that a manipulator running
a hand-written controller does not. The behaviour of a classical controller can be inspected,
bounded and argued about before the machine is switched on; the behaviour of a policy obtained by
reinforcement learning is a consequence of an optimisation run, and is known only through the
statistics of the trajectories it produces. When such a policy is to be deployed on a physical arm
sharing its workspace with people, equipment or the workpiece itself, the question of what it does
in the configurations it visits rarely is not a secondary concern. It is the condition on which
deployment depends at all.

The difficulty is that reinforcement learning, in its standard form, offers a single channel
through which a designer expresses what the agent should and should not do: the reward function.
Safety requirements expressed through that channel become penalty terms, and a penalty term must
be given a weight. That weight is a trade-off rate between task performance and safety which the
designer must fix before knowing what either will be worth, and which has no natural units --- it
is a number tuned until the resulting behaviour looks acceptable. Achiam \emph{et al.}
\cite{achiam2017cpo} make this argument directly in motivating constrained policy optimisation:
for systems that interact physically with or around humans, it is more natural to specify a reward
function \emph{and} a set of constraints than to encode both into one scalar objective.

The alternative is to keep the safety requirement outside the reward and impose it as a constraint
with a budget. A budget has a property the weight does not: it can be stated in advance, in the
units of the quantity being constrained, and it can be audited afterwards against what the trained
policy actually spent. This distinction --- \emph{a shaped reward requires the designer to guess a
weight, a constraint requires only that they state a limit} --- is the premise of this thesis, and
Section~\ref{sec:b-manip} shows that it separates two otherwise similar bodies of work on
singularity avoidance.

This chapter reviews the literature establishing that premise and locates this thesis within it.
Section~\ref{sec:b-saferl} introduces safe reinforcement learning and the taxonomy under which
this work falls. Section~\ref{sec:b-cmdp} gives the constrained Markov decision process, the
formalism used throughout. Section~\ref{sec:b-cpo} covers the deep constrained policy optimisation
algorithms from which the constrained agent is taken. Section~\ref{sec:b-roboticsafety} turns to
safety in robot learning specifically, and Section~\ref{sec:b-manipulation} to reinforcement
learning on manipulators and on the UR5 in particular. Section~\ref{sec:b-manip} treats
manipulability and singularity avoidance, the operative safety quantity of this work.
Section~\ref{sec:b-repro} addresses seed variance. Section~\ref{sec:b-gap} states the gap this
thesis addresses.

\section{Safe reinforcement learning}\label{sec:b-saferl}

Safe reinforcement learning is the study of methods that optimise a policy while respecting some
notion of acceptable behaviour during learning, at deployment, or both. The organising survey of
the field is that of García and Fernández \cite{garcia2015survey}, which divides the literature
into two branches according to where the safety requirement is inserted. The first modifies the
\emph{optimisation criterion}, so that the objective the agent maximises is itself changed --- by
a risk-sensitive transformation, by a worst-case criterion, or by the addition of explicit
constraints. The second modifies the \emph{exploration process}, leaving the objective intact but
restricting or guiding the actions the agent may try, typically using external knowledge such as a
teacher, a demonstration set or a backup controller.

This thesis belongs to the first branch, and to its constrained-criterion sub-family in
particular. The distinction matters for reading Chapter~\ref{ch:results}: methods in the second
branch aim to make the learning process itself safe and are evaluated on violations accumulated
\emph{during} training, whereas methods in the first aim to produce a policy that satisfies a
requirement once trained, and are properly evaluated on the frozen policy. The evaluation protocol
of this work follows the latter convention, and Section~\ref{sec:r-design} explains why an earlier
version of these results, which read safety figures from training-time telemetry, was withdrawn.

The more recent review by Gu \emph{et al.} \cite{gu2024review} surveys the field after the
deep-learning era and reports that the constrained-criterion approach, expressed as a constrained
Markov decision process, has become the prevailing formalism for safe reinforcement learning. That
observation places the present work inside the mainstream of the field rather than at its margin,
and it is why the CMDP is adopted here without further argument.

\begin{kuettable}{Where the safety requirement is inserted: the two branches of safe
reinforcement learning.}{tab:b-taxonomy}
\begin{tabular}{p{3.1cm}p{4.7cm}p{4.4cm}}
\hline
\textbf{Branch} & \textbf{Mechanism} & \textbf{Evaluated on} \\
\hline
Modify the optimisation criterion & Risk-sensitive or worst-case objectives; explicit constraints
on expected cumulative cost & The trained policy, on held-out episodes \\
Modify the exploration process & Teacher guidance, demonstrations, backup controllers, action
shielding & Violations accumulated during training \\
\hline
\end{tabular}
\end{kuettable}

\section{The constrained Markov decision process}\label{sec:b-cmdp}

The formal object underlying the constrained-criterion branch is the constrained Markov decision
process, developed in its modern form by Altman \cite{altman1999cmdp}. A standard Markov decision
process is the tuple (S, A, P, r, \(\gamma\)) over a state space S, an action space A, a transition
kernel P, a reward function r and a discount factor \(\gamma\). A CMDP augments this with a cost
function c and a budget d, and asks for a policy maximising expected return subject to a bound on
expected cumulative cost. Chapter~\ref{ch:methods} states the optimisation problem in full; what
matters here is what the formalism buys.

Two properties are relevant. The first is that the safety requirement is expressed in the units of
the cost rather than in units of reward, so it can be set from a measurement of the system rather
than from a tuning exercise --- the calibration procedure of Section~\ref{sec:m-calib} depends on
this. The second is that the constrained problem admits a Lagrangian relaxation, in which the
constraint is folded into the objective through a multiplier \(\lambda\) that is itself optimised.
The duality theory licensing that treatment is the contribution of \cite{altman1999cmdp}, and it is
what makes the algorithms of the next section possible: rather than solving a constrained problem
directly, they solve a sequence of unconstrained problems in which the penalty weight is not
guessed by the designer but driven by the observed violation.

That last point deserves emphasis, because it is easily lost. A Lagrangian method does end up
optimising a weighted sum of reward and cost, which superficially resembles reward shaping. The
difference is that the weight is not a hyperparameter. It is a dual variable that rises while the
budget is exceeded and falls once it is met, so the designer specifies the budget and the optimiser
discovers whatever weight is required to satisfy it.

\section{Constrained policy optimisation in deep reinforcement learning}\label{sec:b-cpo}

The policy-gradient algorithm used as the baseline throughout this thesis is proximal policy
optimisation \cite{schulman2017ppo}, which limits the size of each policy update through a clipped
surrogate objective and has become the default on-policy method for continuous control. Its
popularity is partly a matter of robustness: it attains much of the reliability of trust-region
methods at a fraction of their implementation complexity, which is also why it is the algorithm
most commonly extended when a constraint is to be added.

Constrained policy optimisation (CPO) \cite{achiam2017cpo} was the first algorithm to bring the
CMDP into deep reinforcement learning at scale. It extends the trust-region approach so that each
update both improves the return and maintains approximate constraint satisfaction, computing a
penalty coefficient afresh at every step. Its significance for this thesis is as much conceptual as
algorithmic: it established the constrained formulation as a practical alternative to reward
shaping for high-dimensional continuous control, and it supplies the motivating argument quoted in
Section~\ref{sec:b-motivation}.

The algorithm actually implemented here, however, is the Lagrangian variant of PPO, as specified in
the Safety Gym benchmark report of Ray \emph{et al.} \cite{ray2019benchmarking}. That report
proposes constrained reinforcement learning as the standard formalism for safe exploration and
supplies a benchmark suite of continuous-control environments together with reference
implementations of PPO, TRPO, their Lagrangian forms, and CPO. Two of its findings bear directly on
the design of this study. The first is the specification itself: the constrained agent studied here
is PPO-Lagrangian as defined in \cite{ray2019benchmarking}, in which an adaptive penalty on the
safety cost is maintained by dual ascent on a Lagrange multiplier. The second is empirical, and is
the reason PPO-Lagrangian was preferred to CPO --- \cite{ray2019benchmarking} reports that CPO
performs poorly on the Safety Gym environments by comparison with the considerably simpler
Lagrangian methods. Given that result, adopting CPO would have required a justification this work
does not have; the Lagrangian form is both easier to implement correctly and better supported by
the published comparison.

Lagrangian methods carry a characteristic weakness, and this thesis encounters it directly in
Chapter~\ref{ch:results}. Because the multiplier is updated by plain dual ascent, its trajectory
oscillates: it rises while the cost exceeds the budget, overshoots, and decays once the cost has
been driven under. Stooke \emph{et al.} \cite{stooke2020pid} analyse this behaviour and propose
treating the multiplier update as a control problem, adding proportional and derivative terms to
damp the oscillation. Their analysis is used in Section~\ref{sec:r-variance} to interpret the
multiplier values observed at the end of training, and the PID form itself is identified in
Chapter~\ref{ch:conclusion} as the natural extension of this work.

The benchmarking landscape has since consolidated. Safety Gymnasium \cite{ji2023safetygymnasium}
supersedes the original Safety Gym environments and ships alongside a library of safe reinforcement
learning algorithms under a common evaluation protocol. Its reporting conventions --- episodic cost
against a stated budget, and violation rates measured per episode --- are those adopted in
Chapter~\ref{ch:results}.

\section{Safety in robot learning}\label{sec:b-roboticsafety}

The literature reviewed so far approaches safety from the machine-learning side. Robotics
approaches it from the control side, with a vocabulary of invariant sets, barrier certificates and
stability guarantees that developed independently. Brunke \emph{et al.} \cite{brunke2022safe}
survey the intersection and, more usefully for the present purpose, reconcile the two vocabularies
into a common framework spanning learning-based control through to safe reinforcement learning.

That reconciliation is necessary here because this thesis sits across the boundary. The quantities
constrained in this work --- the manipulability measure, the distance to a joint limit, the
clearance above the work surface --- are control-theoretic constructs describing the kinematic
state of the arm. The mechanism enforcing them is a learning-theoretic one: a cost critic and a
dual variable. Neither literature alone supplies the language for that combination.

The survey also draws a distinction that constrains what may be claimed from the results of this
work. Methods enforcing a constraint in expectation, as a Lagrangian method does, bound
\emph{expected} exposure to unsafe configurations; they do not bound the severity of any individual
excursion. Methods supplying a safety certificate, such as a control barrier function acting as a
filter on the commanded action, offer the latter guarantee but require a model. This distinction is
returned to in Section~\ref{sec:r-safety}, where a counter-result of exactly this shape is
reported, and in Chapter~\ref{ch:conclusion}, where the two mechanisms are argued to address
different halves of the problem.

Within manufacturing specifically, Elguea-Aguinaco \emph{et al.} \cite{elguea2023review} review
reinforcement learning for contact-rich manipulation from 2017 onward, covering reward design,
sim-to-real transfer and the treatment of safety. Their survey is used in Chapter~\ref{ch:methods}
to establish which of the design choices made in this work follow established practice.

\section{Deep reinforcement learning for robotic manipulation}\label{sec:b-manipulation}

Applications of deep reinforcement learning to manipulation divide broadly by algorithm family.
Shahid \emph{et al.} \cite{shahid2022continuous_grasping} compare an on-policy method (PPO) against
an off-policy method (soft actor-critic) on a grasping task with a Franka Emika Panda, using a
dense reward and a fine-tuning procedure that adapts a trained policy to a modified task. The study
is relevant here in three ways: it is the methodological template for an on-policy versus
off-policy comparison on a manipulation task; it demonstrates transfer of the learned behaviour
from simulation to a physical robot, which supports the position taken in
Chapter~\ref{ch:conclusion} that results of this kind are transferable in principle; and it is the
precedent for the off-policy arm that was planned for this study and ultimately not trained, as
recorded in Section~\ref{sec:r-limits}.

Work specific to the UR5 platform is less common. Ferreira and Barbosa \cite{ferreira2025grasping}
train a UR5 fitted with a Robotiq 2F-85 gripper to grasp randomly positioned objects in a PyBullet
environment, comparing soft actor-critic against PPO under a reward combining distance penalties
with grasp and pose terms. They report grasp success rates of 87\,\% for SAC and 82\,\% for PPO.
These are the closest published figures for a learned UR5 grasping policy, and they are quoted here
with one qualification that must be carried into any comparison: their task requires a full
physical grasp, in which the fingers close on the object and the resulting contact must be stable,
whereas the task studied in this thesis abstracts the grasp as a weld for the reasons given in
Section~\ref{sec:m-env}. The success rates of Chapter~\ref{ch:results} are therefore not directly
comparable to theirs, and are not presented as though they were.

On the safety side, Xia \emph{et al.} \cite{xia2024proactive} apply deep reinforcement learning to
proactive obstacle avoidance for a UR5e in a human-robot collaborative assembly setting, with depth
sensing for environmental perception and a wrist force-torque sensor. Their work establishes that
safe operation of this specific manipulator is an active industrial concern rather than an academic
construction, which is the premise on which the practical relevance of this thesis rests.

\section{Manipulability and singularity avoidance}\label{sec:b-manip}

The operative safety constraint in this thesis is a floor on manipulability. The measure is due to
Yoshikawa \cite{yoshikawa1985manipulability}, who proposed \(w = \sqrt{\det(JJ^{\top})}\), computed
from the manipulator Jacobian J, as a scalar index of a configuration's capacity to generate
end-effector motion in arbitrary directions. The index falls towards zero as the arm approaches a
kinematic singularity, at which it loses the ability to move instantaneously along some Cartesian
direction and at which small commanded end-effector velocities map to large joint velocities. It is
the standard measure of this property and is used throughout this work; Chapter~\ref{ch:methods}
defines the cost term built from it.

Singularity avoidance has been treated within learning frameworks before, and the manner of that
treatment is precisely where this thesis differs. Shen \emph{et al.} \cite{shen2022reactive}
address reactive obstacle avoidance for redundant manipulators by learning motion in the null space
of the Jacobian, and introduce the manipulability index \emph{into the reward function} so that the
arm maintains a well-conditioned posture while executing its task. The approach is sound and the
objective is the same as the one pursued here. The mechanism is not. Placing manipulability in the
reward requires a weight relating a unit of manipulability to a unit of task progress, and that
weight is a design choice with no principled setting. Placing it in a constraint, as this thesis
does, requires instead a budget --- a statement of how much cumulative proximity to singularity an
episode may incur --- which can be measured from a baseline policy and audited against the trained
one.

The comparison is worth stating plainly, because it is the sharpest available expression of what
this thesis does differently. Both approaches seek the same behaviour. One obtains it by tuning a
trade-off; the other by declaring a limit.

\section{Reproducibility and seed variance}\label{sec:b-repro}

A final strand of the literature is methodological rather than algorithmic, and it bears on the
central experimental finding of this work. Henderson \emph{et al.} \cite{henderson2018matters}
examine the reproducibility of results in deep reinforcement learning and demonstrate that outcomes
for a single algorithm on a single task can differ substantially according to the random seed
alone, independently of any change to the algorithm, its hyperparameters or the environment. Their
recommendation is that results be reported over multiple independent seeds with a measure of
dispersion, and that comparisons between methods be treated with corresponding caution.

Two consequences follow for this thesis. The first is procedural: the protocol of
Chapter~\ref{ch:methods} trains every arm on ten independent seeds and reports dispersion across
them rather than a single figure, and that protocol is justified by reference to
\cite{henderson2018matters} rather than by assertion. The second is interpretive: the seed-to-seed
spread reported in Section~\ref{sec:r-variance} is not an anomaly of this environment but an
instance of a documented property of the method, which is why it is presented there as a measured
quantity rather than as a surprising discovery.

\section{Research gap and positioning}\label{sec:b-gap}

The nearest published analogue to the present study is that of Khan \emph{et al.}
\cite{khan2026rl_precision_grasping}, which compares a model-free PPO baseline against its
constrained variant for precision grasping and safety-critical coordination on a robotic arm, using
the Safety Gym framework extended with a Panda manipulator, and introduces a tactile feedback
scheme to accelerate learning under safety compliance. It is the closest match in method, in
framing and in intent, and this thesis should be read as continuing that line of work rather than
as departing from it.

Two gaps are nonetheless visible in that line, and it is these the present study addresses.

The first is a matter of experimental control. A constrained agent of this family differs from its
unconstrained baseline in more than the constraint: it constructs and trains an additional cost
critic, whose presence can perturb the optimiser through shared machinery even when the constraint
itself is inactive. Unless a third arm is trained that carries the cost critic but pins the
multiplier to zero, any measured difference between the constrained agent and the baseline sums two
effects of different kinds, and cannot be attributed to the constraint alone. Published comparisons
in this area, including \cite{khan2026rl_precision_grasping}, do not report such a control arm.
Chapter~\ref{ch:results} shows this is not a hypothetical concern: an earlier version of the
present work reported a large advantage for the constrained agent that was subsequently traced to
exactly this class of implementation artifact and withdrawn in full.

The second is a matter of dispersion. Comparisons in this literature are typically reported as
means, without an account of how much the outcome varies across training runs differing only in
their random seed. Given the finding of \cite{henderson2018matters} discussed in
Section~\ref{sec:b-repro}, a mean taken over few seeds is a weak basis for a safety claim: an
integrator deploying such a policy does not obtain the mean over training runs they will never
perform, but the single policy they happened to train. Whether a safety constraint changes the
\emph{predictability} of the outcome, and not merely its average, is therefore a question of
practical consequence, and it is not one the existing literature answers.

This thesis addresses both. It trains a three-arm comparison --- an unconstrained baseline, a
control arm carrying the cost critic with the multiplier pinned to zero, and the constrained agent
--- on ten independent seeds each, and reports the decomposition of the constrained-versus-baseline
difference into an implementation term and a constraint term, together with the dispersion of every
reported quantity across seeds.

The work is also positioned locally. It continues a line of manipulator research in the same
laboratory, of which the immediately preceding study \cite{khan2025csrt_ibvs} implemented classical
image-based visual servoing with CSRT tracking on a five-degree-of-freedom manipulator. Where that
work addressed the perception and servoing loop with classical methods, the present work addresses
the control policy with learned ones, and adds the safety constraint that classical visual servoing
does not itself provide.

\section{Summary}\label{sec:b-summary}

Safety requirements on a learned manipulation policy are poorly served by reward shaping, which
obliges the designer to fix a trade-off weight in advance, and are better expressed as constraints
with stated budgets \cite{achiam2017cpo}. The constrained Markov decision process
\cite{altman1999cmdp} is the formalism for this, and is the prevailing one in contemporary safe
reinforcement learning \cite{gu2024review}. Among the deep algorithms implementing it, the
Lagrangian variant of PPO \cite{schulman2017ppo, ray2019benchmarking} is both simpler than the
trust-region alternative \cite{achiam2017cpo} and better supported by the published benchmark
comparison, and is the algorithm adopted here.

Applied to manipulators, this machinery has been used for obstacle avoidance and for grasping, on
the UR5 among other platforms \cite{ferreira2025grasping, xia2024proactive}, and singularity
avoidance in particular has been pursued by placing the manipulability measure
\cite{yoshikawa1985manipulability} inside the reward \cite{shen2022reactive}. This thesis places it
inside a constraint instead. Two gaps in the nearest published comparison
\cite{khan2026rl_precision_grasping} define the contribution: the absence of a control arm
isolating the effect of the cost critic from the effect of the constraint, and the absence of any
account of seed-to-seed dispersion. Chapter~\ref{ch:methods} sets out the method by which both are
addressed.
